\begin{recipe}{Schiacciata}
\kicker[Brood,Toscaans,Basisrecept]{Brood \& basis · Toscaans · basisrecept}
\index[register]{Schiacciata}
\index[register]{Bloem}
\index[register]{Rozemarijn}
\index[register]{Gist}

\meta{Rijzen ca.\ 1,5 uur \dvd 1 bakplaat \dvd Oven 220\,°C}
\marginimage{images/schiacciata-ingredienten}

\lettrine{D}{eze} schiacciata aten we voor het eerst bij een panini caprese op
de Passo Albetone in Italië (zie blz.~\pageref{recipe:panini-caprese}), en we
dachten eerst dat het gewoon focaccia was. Schiacciata betekent letterlijk
platgedrukt en lijkt qua bereiding ook op focaccia, maar is dunner en
knapperiger in plaats van dik en luchtig. Daardoor snijd je het makkelijk
horizontaal open om te beleggen met vleeswaren en kaas.

\blockrule{Ingrediënten}
\begin{ingredients}
  \ing{500 g}{bloem tipo 00}\index[register]{Bloem}
  \ing{350 g}{lauwwarm water (70\% hydratatie)}
  \ing{20 g}{extra vergine olijfolie (ca.\ 2 el, 4\%)}
  \ing{10 g}{zout (2\%)}
  \ing{7 g}{droge gist (1 zakje, 1,4\%)}
  \ingb{extra vergine olijfolie, om te besprenkelen}
  \ingb{grof zeezout}
  \ingb{verse rozemarijn, optioneel}\index[register]{Rozemarijn}
\end{ingredients}

\blockrule{Bereiding}
\tip{Het deeg is relatief nat en plakkerig: dat hoort zo en is precies wat
schiacciata zijn luchtige binnenkant en knapperige korst geeft. Wil je hem
gebruiken als broodje voor beleg, bijvoorbeeld voor een panini caprese (zie
blz.~\pageref{recipe:panini-caprese})? Laat de rozemarijn dan achterwege,
zonder is hij veelzijdiger inzetbaar.}

\medskip
\noindent\textit{Deeg maken}
\begin{steps}
  \item Meng de bloem en het zout in een grote mengkom.
  \item Los de gist op in het lauwwarme water.
  \item Giet het gistwater en de olijfolie bij de bloem.
  \item Kneed het deeg ongeveer 10 minuten tot het soepel en elastisch is.
\end{steps}

\medskip
\noindent\textit{Eerste rijs}
\begin{steps}
  \item Vet een schone kom lichtjes in met olijfolie. Leg de deegbal erin en
        dek de kom af met een vochtige theedoek. Laat het deeg 1 uur rijzen op
        een warme plek tot het in volume is verdubbeld.
\end{steps}

\medskip
\noindent\textit{Vormen op de bakplaat}
\begin{steps}
  \item Vet een grote bakplaat rijkelijk in met olijfolie (of gebruik
        bakpapier). Stort het deeg op de plaat.
  \item Duw het deeg met je vingertoppen voorzichtig uit tot een gelijkmatige,
        rechthoekige lap. Dek weer af en laat nog 30 minuten rusten.
\end{steps}

\medskip
\noindent\textit{Platdrukken en opmaken}
\begin{steps}
  \item Verwarm de oven ondertussen voor op 220\,°C (boven- en onderwarmte).
  \item Druk met je vingertoppen diepe kuiltjes in het gerezen deeg, tot aan
        de bodem.
  \item Besprenkel de bovenkant gul met extra olijfolie, zodat de kuiltjes
        volstaan. Strooi het grove zeezout en de rozemarijn gelijkmatig over
        het deeg.
\end{steps}

\medskip
\noindent\textit{Bakken}
\begin{steps}
  \item Schuif de bakplaat in het midden van de voorverwarmde oven. Bak de
        schiacciata in 20 tot 25 minuten goudbruin en knapperig.
  \item Laat na het bakken kort afkoelen op een rooster.
\end{steps}

\heroimagefade{images/schiacciata-hero}
\end{recipe}
